% Standalone section -- paste into the weekly report, or \input it.
% Week 3 working note: whether the tortuosity panel (tortuosity_and_representation.tex,
% topology/tortuosity.py) separates healthy from diseased ImageCAS-X cases, and whether the
% metrics scripts/tortuosity_sensitivity.py found most reproducible under a smoothing-scale sweep
% are the ones that do. Companion to bifurcation_angle_and_disease.tex, which asks the same
% disease-separation question of bifurcation angle and finds a well-powered null; this one finds
% the opposite shape of result, a real signal concentrated in the metrics least trusted for
% reproducibility. Numbers from scripts/tortuosity_disease.py, run 2026-09-16 (800 cases).

\section*{Tortuosity and disease: the reproducible metric is not the discriminating one}

The smoothing-scale sensitivity study that followed the tortuosity panel's first description
settled which of its metrics survive a defensible range of preprocessing choices: non-planarity
and the out-of-plane distance family, not curvature or the turning-angle sum, which reorder the
cohort under a different sigma badly enough to reshuffle a top-5\% tail
(tortuosity\_and\_representation.tex). The natural next question is whether those same reproducible
metrics are the ones that separate ImageCAS-X's diseased cases from its healthy ones, on the same
800 cases the bifurcation angle comparison already tested (bifurcation\_angle\_and\_disease.tex).
They are not. Eighteen of sixty vessel-by-metric tests survive Benjamini--Hochberg correction at
5\% (\texttt{scripts/tortuosity\_disease.py}), and the largest, most significant effects sit
entirely in curvature and the turning-angle sum, the metrics flagged least stable under a
smoothing-scale change, while non-planarity and the out-of-plane distance family, recommended for
exactly this kind of extreme-value flag, come back null in every vessel. The one metric that is
both reproducible and disease-associated is the plain distance ratio, $L/D$, on the LAD.

\subsection*{Curvature and turning rate separate the cohorts; the metrics kept for reproducibility do not}

The RCA and LCx show the clearest effect. Diseased RCAs turn 12.7\% more per millimeter than
healthy ones (mean turning-angle sum 0.0742 to 0.0836 rad/mm, $n_h=412$, $n_d=386$,
$p=4.4\times10^{-13}$) and their mean curvature is 12.3\% higher (0.0743 to 0.0834 mm$^{-1}$,
$p=1.3\times10^{-12}$); the LCx shows the same pattern at similar magnitude (turning-angle sum
$+12.5\%$, $p=1.1\times10^{-7}$; mean curvature $+12.2\%$, $p=1.9\times10^{-7}$). The LAD shows the
same direction at roughly half the effect size (mean curvature $+6.7\%$, $p=0.0002$). Even torsion,
the noisiest metric in the whole panel, comes back significant on the RCA ($+4.5\%$,
$p=3.0\times10^{-6}$), though its effect is the smallest of any significant result, consistent with
a real but small signal sitting inside a mostly noise-dominated measurement. Against that,
non-planarity and the out-of-plane distance family are flat in every vessel: RCA non-planarity
moves $+20.4\%$ in the same direction but does not clear $p=0.15$, and the other three vessels move
by less and no more significantly.

One metric in the panel matches the clinical literature's own definition directly,
three or more bends exceeding $45^\circ$ \cite{li2011tortuosity, zebicmihic2023tortuosity}, and it
is null everywhere: the median bend count is identical between healthy and diseased cases in all
four vessels (LM 0, LAD 4, LCx 3, RCA 3), and no vessel's difference clears $p=0.07$. The clinical
bend count and the continuous curvature measures are answering different questions on the same
centerlines, and only the continuous ones separate this cohort.

\subsection*{Not an artifact of image quality}

Diseased scans in ImageCAS-X carry a slightly lower quality grade than healthy ones (median 4 vs.\
3 on the dataset's 4-point scale, $n=800$, $p=0.003$), which is exactly the confound that would
manufacture a spurious curvature effect: noisier segmentations produce noisier centerlines, and
curvature is the metric most sensitive to that noise. Restricting the comparison to
Image Quality~$=4$ only (the best grade, $n_h\approx247$, $n_d\approx187$--$189$ depending on
vessel) removes the confound rather than assuming it away, and the RCA and LCx effects survive
essentially unchanged: RCA mean curvature $+10.6\%$ ($p=4.8\times10^{-6}$), RCA turning-angle sum
$+9.7\%$ ($p=4.2\times10^{-6}$), LCx turning-angle sum $+10.2\%$ ($p=0.002$). The LAD's curvature
effect weakens to borderline within this stratum ($p=0.02$--$0.03$), but its $L/D$ effect does not
($+5.2\%$, $p=0.003$). Whatever is driving the RCA and LCx result, it is not simply that diseased
scans are worse images.

\subsection*{The null on the reproducible metrics is well powered for the LAD, not for the others}

A null result needs its own power check before it licenses anything, and the answer here differs
by vessel. On the LAD, this cohort's size gives 80\% power to detect an effect as small as $12.9\%$
of the healthy median for non-planarity, $5.6\%$ for the out-of-plane distance ratio, and $6.0\%$
for the out-of-plane angle: an effect of the size seen in curvature ($+6.7\%$) would very likely
have been caught if it were present in these metrics too, so the LAD's null is a genuine one, not a
power shortfall. On the RCA, the out-of-plane distance and angle nulls are similarly well powered
(minimum detectable effect $11$--$12\%$ of the healthy median), but the non-planarity null is not
($49\%$): RCA vessels are close enough to planar, and non-planarity's own spread across the cohort
small enough, that only a very large relative shift would register at all. The LM is uninformative
on every reproducible metric: its plane fit is degenerate on most cases (its own best-fit plane is
undefined on a vessel this short and this close to straight, topology/tortuosity.py), which drops
its usable $n$ from $\sim\!400$ to $82$ vs.\ $85$ and pushes the minimum detectable effect past
$33\%$ everywhere. An LM tortuosity--disease null, on this evidence, says nothing either way.

\subsection*{The direction does not resolve against either published pattern}

The largest tortuosity cohort found, Groves et al.'s 1,221 consecutive angiograms, associates
severe tortuosity with a significantly \emph{lower} incidence of significant coronary artery
disease ($p=0.003$) \cite{groves2009tortuosity}; Li et al., in 1,010 angiography patients using the
same three-bends-over-$45^\circ$ definition, found the same inverse direction (odds ratio $0.755$,
95\% CI $0.574$--$0.994$, $p=0.045$) and no difference in major adverse events over 2--4 years
\cite{li2011tortuosity}. Zebić Mihić et al., in 131 patients, found the opposite sign but for a
different endpoint: tortuosity was associated with ischemia specifically in the \emph{absence} of
obstructive disease (odds ratio $7.96$) \cite{zebicmihic2023tortuosity}, and a second, 160-patient
cohort from the same group showed a continuous tortuosity index detects territory-matched ischemia
associations that the clinical bend count misses almost entirely
\cite{zebicmihic2024index}. ImageCAS-X's \texttt{Disease} field is a single yes/no with no severity
grade and no obstructive/non-obstructive distinction, so this comparison cannot say which of those
two literatures, if either, the RCA and LCx effect agrees with: a coarse label large enough to
contain both obstructive and non-obstructive disease could produce a positive average of two
populations that individually run in opposite directions. What it does establish is that "more
curved" is not obviously the direction the obstructive-disease literature would predict, which is
reason enough not to read this result as confirming the field's more common assumption that
tortuosity and disease simply rise together.

\subsection*{What this licenses}

Within ImageCAS-X, the RCA and LCx bend measurably more in cases flagged diseased, an effect that
survives multiple-comparison correction and an image-quality stratification built specifically to
rule out the most obvious confound; the LAD's $L/D$ carries the same signal at smaller size while
also being the one metric in the panel already shown stable under a smoothing-scale sweep. What
this does not establish is which literature that effect agrees with, since the two published
directions require a severity grade this dataset's label cannot provide, and it does not establish
that non-planarity or the out-of-plane family are uninformative in general, only that they are
uninformative against this particular coarse label, on vessels short enough that several of the
tests were never well powered to see anything else. \texttt{Disease} is also a per-patient flag
tested against a per-vessel measurement, exactly the dilution
bifurcation\_angle\_and\_disease.tex's own literature comparison already named as the gap between a
site-matched clinical design and a whole-tree present/absent label, and there is still no age, sex
or BMI field in ImageCAS-X to check against, the same confounder gap that comparison could only flag rather than close.
A severity-graded, site-matched redesign is where this belongs next, and the project plan already
names it: the CGPS baseline-vs-progression cohort, once it exists, is where tortuosity's
association with disease should be tested for real, not reproduced from a two-value label a second
time.
